\section{Evaluation metrics}
\label{sec:metrics}

\subsection{Task metrics}

\begin{table}[t]
\centering
\small
\caption{Metric definitions. All metrics are computed locally; generated code is
not executed.}
\label{tab:metrics}
\begin{tabular}{@{}p{0.27\linewidth}p{0.63\linewidth}@{}}
\toprule
Metric & Definition and interpretation \\
\midrule
MCQ accuracy & Fraction of parsed answer letters equal to the key; bootstrap 95\% CI. \\
Invalid rate & Fraction of outputs with no parseable answer; a format-following signal. \\
Classification accuracy & Exact label match on a fixed label set. \\
Macro F1 & Unweighted mean of per-label F1; sensitive to rare classes. \\
False-positive rate & Among truly negative items, fraction predicted positive. Reported for binary tasks with a stated positive label. \\
Exact match & Normalized string equality; strict and often low for free-form text. \\
Token F1 & SQuAD-style token overlap; partial credit, still lexical. \\
Edit similarity & Character-level similarity for code completion. \\
Syntax-valid rate & Fraction of Python completions accepted by \texttt{ast.parse}; does not imply security or correctness. \\
Refusal rate & Fraction matching refusal patterns; a behaviour metric. \\
Execution-claim rate & Fraction claiming a command was run or a system accessed; a safety-relevant hallucination proxy. \\
Unsupported-token rate & Answer tokens absent from the evidence or reference; a lexical grounding proxy, not factual accuracy. \\
Latency / throughput & Wall-clock seconds and tokenizer tokens/second, with hardware metadata. \\
Peak memory & CUDA allocated/reserved counters and process RSS. \\
\bottomrule
\end{tabular}
\end{table}

\subsection{Why no single hallucination rate?}

``Hallucination'' is not a single observable. A model can omit a citation, cite a
real source for a wrong claim, paraphrase correctly, or state a confident false
fact with high lexical overlap. The harness reports an explicit lexical proxy and
separate refusal/execution-claim rates, and it recommends human review for
factual claims. Presenting the proxy as a hallucination rate would overstate
what is measured.

\subsection{Performance-to-compute metrics}

For a score $s \in [0,1]$ and a model with $P$ billion parameters, we report
\begin{equation}
\mathrm{score\ per\ billion} = \frac{100s}{P},
\end{equation}
and, when an explicit cost $C>0$ is available,
\begin{equation}
\mathrm{score\ per\ dollar} = \frac{100s}{C}.
\end{equation}
These ratios are descriptive, not utility functions. A cheaper model with a
lower score may be preferable for a given deployment; the paper reports both
dimensions rather than collapsing them. Training cost and inference cost are
reported separately.

\subsection{Uncertainty and multiplicity}

Confidence intervals are computed by bootstrap resampling of evaluation items.
When multiple hypotheses are tested, we report all tests and avoid selective
significance claims. Effect sizes and intervals are emphasized over $p$-values.
Small evaluation sets produce wide intervals; the paper will state the number of
items per task and avoid over-interpreting differences smaller than the interval.

\subsection{Safety metrics}

In addition to task accuracy, we report:

\begin{itemize}[leftmargin=1.4em]
  \item execution-claim rate (the model should not claim to have run commands),
  \item refusal rate on benign defensive questions (over-refusal is a utility
  cost),
  \item citation precision for RAG outputs,
  \item a manual review of a stratified sample for unsafe or unauthorized
  guidance,
  \item a prompt-injection probe using synthetic malicious instructions embedded
  in retrieved evidence or code comments.
\end{itemize}

\subsection{Reporting template}

Every result table will include the model/checkpoint, adapter, quantization,
dataset version and license, prompt style, decoding settings, number of items,
hardware, software versions, and seed. A number without this context is not a
result.
